\documentclass{Phyport}

\usepackage{hyperref}
\hypersetup{
    colorlinks=true,
}
% % 选用宏（如有额外需要的宏包请在此添加）
\usepackage{float}       % 提供 [H] 选项
\usepackage{booktabs}    % 三线表
\usepackage{amsmath}     % 数学公式
\usepackage{caption}     % 标题设置
\usepackage{siunitx}     % SI 单位排版
 %选用宏（如有）

% ========== 额外需要的宏包 ==========
\usepackage{float}       % 提供 [H] 选项，强制浮动体在此处
\usepackage{booktabs}    % 三线表必备（\toprule, \midrule, \bottomrule）
\usepackage{multirow}    % 合并行
\usepackage{makecell}    % 单元格内换行（\makecell）
\usepackage{caption}     % 更灵活的标题设置
\usepackage{subcaption}  % 子图
\usepackage{amsmath}     % 数学公式
\usepackage{siunitx}     % SI 单位排版
\usepackage{tabularx}    % 自动列宽表格

% 信息（示例用）
\exname{功能实例演示}
\extable{00}
\instructor{张三}
\class{物理2501}
\name{李四}
\stuid{12345678}

\nyear{2026}
\nmonth{3}
\nday{14}
\nweekday{六}

\redate{}
\resitu{}

\begin{document}
\setcounter{page}{0}
\makecover

% ============================================================
%                     一、图片相关
% ============================================================
\section{图片插入示例}

% -------------------- 1.1 单张图片 --------------------
\subsection{单张图片}

使用 \verb|\includegraphics| 插入单张图片，搭配 \texttt{figure} 环境和 \texttt{[H]} 选项固定位置。

\begin{figure}[H]
    \centering
    \includegraphics[width=0.4\linewidth]{example-image} % LaTeX 内置测试图
    \caption{单张图片示例}
    \label{fig:single}
\end{figure}

引用方式：如图~\ref{fig:single} 所示。

% -------------------- 1.2 指定尺寸 --------------------
\subsection{不同尺寸控制方式}

\begin{figure}[H]
    \centering
    % 方式1：按行宽比例
    \includegraphics[width=0.3\linewidth]{example-image-a}
    \caption{宽度 = 0.3 倍行宽}
    \label{fig:size1}
\end{figure}

\begin{figure}[H]
    \centering
    % 方式2：指定绝对尺寸
    \includegraphics[width=6cm, height=3cm]{example-image-b}
    \caption{指定宽6cm、高3cm（可能变形）}
    \label{fig:size2}
\end{figure}

\begin{figure}[H]
    \centering
    % 方式3：保持比例限制最大尺寸
    \includegraphics[width=6cm, height=3cm, keepaspectratio]{example-image-c}
    \caption{宽6cm、高3cm、保持纵横比（不变形）}
    \label{fig:size3}
\end{figure}

% -------------------- 1.3 并排双图 --------------------
\subsection{并排双图}

使用 \texttt{subcaption} 宏包的 \texttt{subfigure} 环境实现并排子图。

\begin{figure}[H]
    \centering
    \begin{subfigure}[b]{0.45\linewidth}
        \centering
        \includegraphics[width=\linewidth]{example-image-a}
        \caption{子图A}
        \label{fig:sub-a}
    \end{subfigure}
    \hfill
    \begin{subfigure}[b]{0.45\linewidth}
        \centering
        \includegraphics[width=\linewidth]{example-image-b}
        \caption{子图B}
        \label{fig:sub-b}
    \end{subfigure}
    \caption{并排双图示例}
    \label{fig:side-by-side}
\end{figure}

% -------------------- 1.4 三图并排 --------------------
\subsection{三图并排}

\begin{figure}[H]
    \centering
    \begin{subfigure}[b]{0.3\linewidth}
        \centering
        \includegraphics[width=\linewidth]{example-image-a}
        \caption{子图1}
    \end{subfigure}
    \hfill
    \begin{subfigure}[b]{0.3\linewidth}
        \centering
        \includegraphics[width=\linewidth]{example-image-b}
        \caption{子图2}
    \end{subfigure}
    \hfill
    \begin{subfigure}[b]{0.3\linewidth}
        \centering
        \includegraphics[width=\linewidth]{example-image-c}
        \caption{子图3}
    \end{subfigure}
    \caption{三图并排示例}
    \label{fig:three-side}
\end{figure}

% -------------------- 1.5 图片旋转与裁剪 --------------------
\subsection{图片旋转与裁剪}

\begin{figure}[H]
    \centering
    \includegraphics[width=0.3\linewidth, angle=45]{example-image}
    \caption{旋转45°的图片}
    \label{fig:rotate}
\end{figure}

\begin{figure}[H]
    \centering
    % trim={left bottom right top}, clip 必须配合使用
    \includegraphics[width=0.4\linewidth, trim={1cm 1cm 1cm 1cm}, clip]{example-image}
    \caption{四周各裁剪1cm的图片}
    \label{fig:trim}
\end{figure}


% ============================================================
%                     二、表格相关（三线表）
% ============================================================
\section{表格插入示例（三线表）}

% -------------------- 2.1 基础三线表 --------------------
\subsection{基础三线表}

三线表使用 \texttt{booktabs} 宏包提供的 \verb|\toprule|、\verb|\midrule|、\verb|\bottomrule|。

\begin{table}[H]
    \centering
    \caption{测量数据记录表}
    \label{tab:basic}
    \begin{tabular}{ccc}
        \toprule
        序号 & 测量值 $x$ / mm & 偏差 $\Delta x$ / mm \\
        \midrule
        1 & 25.32 & 0.02 \\
        2 & 25.28 & $-0.02$ \\
        3 & 25.30 & 0.00 \\
        4 & 25.34 & 0.04 \\
        5 & 25.26 & $-0.04$ \\
        \bottomrule
    \end{tabular}
\end{table}

引用方式：如表~\ref{tab:basic} 所示。

% -------------------- 2.2 带单位的三线表 --------------------
\subsection{带单位的三线表（siunitx）}

使用 \texttt{siunitx} 宏包的 \texttt{S} 列类型，可以自动对齐小数点。

\begin{table}[H]
    \centering
    \caption{不同温度下的电阻值}
    \label{tab:siunitx}
    \begin{tabular}{c S[table-format=3.1] S[table-format=4.2]}
        \toprule
        序号 & {温度 $T$ / \si{\degreeCelsius}} & {电阻 $R$ / \si{\ohm}} \\
        \midrule
        1 &  20.0 &  100.23 \\
        2 &  40.0 &  115.47 \\
        3 &  60.0 &  131.89 \\
        4 &  80.0 &  148.56 \\
        5 & 100.0 &  165.12 \\
        \bottomrule
    \end{tabular}
\end{table}

% -------------------- 2.3 合并列 --------------------
\subsection{合并列的三线表}

\begin{table}[H]
    \centering
    \caption{多次测量汇总}
    \label{tab:multicolumn}
    \begin{tabular}{cccc}
        \toprule
        \multicolumn{2}{c}{正向测量} & \multicolumn{2}{c}{反向测量} \\
        \cmidrule(lr){1-2} \cmidrule(lr){3-4}
        $U$ / V & $I$ / mA & $U$ / V & $I$ / mA \\
        \midrule
        1.0 & 5.2  & 1.0 & 5.1 \\
        2.0 & 10.3 & 2.0 & 10.1 \\
        3.0 & 15.1 & 3.0 & 15.3 \\
        4.0 & 20.4 & 4.0 & 20.2 \\
        \bottomrule
    \end{tabular}
\end{table}

% -------------------- 2.4 合并行 --------------------
\subsection{合并行的三线表}

\begin{table}[H]
    \centering
    \caption{分组实验数据}
    \label{tab:multirow}
    \begin{tabular}{cccc}
        \toprule
        材料 & 序号 & 长度 / cm & 质量 / g \\
        \midrule
        \multirow{3}{*}{铜} & 1 & 10.0 & 89.6 \\
                            & 2 & 15.0 & 134.4 \\
                            & 3 & 20.0 & 179.2 \\
        \midrule
        \multirow{3}{*}{铝} & 1 & 10.0 & 27.0 \\
                            & 2 & 15.0 & 40.5 \\
                            & 3 & 20.0 & 54.0 \\
        \bottomrule
    \end{tabular}
\end{table}

% -------------------- 2.5 自动列宽表格 --------------------
\subsection{自动列宽三线表（tabularx）}

\texttt{tabularx} 的 \texttt{X} 列类型会自动分配剩余列宽，适合长文字内容。

\begin{table}[H]
    \centering
    \caption{实验仪器清单}
    \label{tab:tabularx}
    \begin{tabularx}{\linewidth}{c X c}
        \toprule
        序号 & 仪器名称及规格 & 数量 \\
        \midrule
        1 & 游标卡尺（量程 0--150\,mm，分度值 0.02\,mm） & 1 \\
        2 & 螺旋测微器（量程 0--25\,mm，分度值 0.01\,mm） & 1 \\
        3 & 电子天平（量程 0--200\,g，感量 0.01\,g） & 1 \\
        4 & 待测样品（铜柱、铝柱若干） & 若干 \\
        \bottomrule
    \end{tabularx}
\end{table}

% -------------------- 2.6 单元格内换行 --------------------
\subsection{单元格内换行（makecell）}

\begin{table}[H]
    \centering
    \caption{含多行表头的数据表}
    \label{tab:makecell}
    \begin{tabular}{ccc}
        \toprule
        序号 & \makecell{测量值 \\ $x$ / mm} & \makecell{不确定度 \\ $u(x)$ / mm} \\
        \midrule
        1 & 25.32 & 0.03 \\
        2 & 25.28 & 0.03 \\
        3 & 25.30 & 0.03 \\
        \bottomrule
    \end{tabular}
\end{table}


% ============================================================
%               三、公式相关
% ============================================================
\section{公式示例}

% -------------------- 3.1 行内公式 --------------------
\subsection{行内公式}

牛顿第二定律 $F = ma$，其中 $F$ 为合力，$m$ 为质量，$a$ 为加速度。

% -------------------- 3.2 行间公式（带编号） --------------------
\subsection{行间公式（带编号）}

\begin{equation}
    E = mc^2
    \label{eq:einstein}
\end{equation}

如公式~\eqref{eq:einstein} 所示。

% -------------------- 3.3 多行公式对齐 --------------------
\subsection{多行公式对齐}

\begin{align}
    \bar{x} &= \frac{1}{n} \sum_{i=1}^{n} x_i \label{eq:mean} \\
    s &= \sqrt{\frac{1}{n-1} \sum_{i=1}^{n} (x_i - \bar{x})^2} \label{eq:std} \\
    u_A &= \frac{s}{\sqrt{n}} \label{eq:ua}
\end{align}

其中公式~\eqref{eq:mean} 为平均值，公式~\eqref{eq:std} 为标准差，公式~\eqref{eq:ua} 为 A 类不确定度。

% -------------------- 3.4 无编号公式 --------------------
\subsection{无编号公式}

\begin{equation*}
    \Delta = \frac{|x_{\text{测}} - x_{\text{真}}|}{x_{\text{真}}} \times 100\%
\end{equation*}

% -------------------- 3.5 带单位的结果表达 --------------------
\subsection{带单位的结果表达（siunitx）}

测量结果为 $L = \SI{25.30 \pm 0.03}{mm}$。

温度 $T = \SI{20.0}{\degreeCelsius}$，电阻 $R = \SI{100.23}{\ohm}$。


% ============================================================
%                 四、其他常用功能
% ============================================================
\section{其他常用功能}

% -------------------- 4.1 有序列表 --------------------
\subsection{有序列表}

\begin{enumerate}
    \item 调节仪器水平
    \item 校准零点
    \item 进行正向测量
    \item 进行反向测量
    \item 记录数据
\end{enumerate}

% -------------------- 4.2 无序列表 --------------------
\subsection{无序列表}

\begin{itemize}
    \item 游标卡尺
    \item 螺旋测微器
    \item 电子天平
\end{itemize}

% -------------------- 4.3 代码插入 --------------------
\subsection{代码插入}

模版中 \texttt{listings} 已配置好，可直接使用。

\begin{lstlisting}[language=Python, caption={Python 数据处理示例}]
import numpy as np

data = np.array([25.32, 25.28, 25.30, 25.34, 25.26])
mean = np.mean(data)
std = np.std(data, ddof=1)
print(f"mean = {mean:.2f}, std = {std:.4f}")
\end{lstlisting}

% -------------------- 4.4 脚注 --------------------
\subsection{脚注}

仪器精度等级为 0.5 级\footnote{精度等级表示仪器的最大引用误差百分数。}。

% -------------------- 4.5 引用汇总 --------------------
\subsection{交叉引用汇总}

\begin{itemize}
    \item 引用图片：图~\ref{fig:single}
    \item 引用表格：表~\ref{tab:basic}
    \item 引用公式：公式~\eqref{eq:einstein}
    \item 引用子图：图~\ref{fig:sub-a} 和图~\ref{fig:sub-b}
\end{itemize}

% -------------------- 4.6 超链接 --------------------
\subsection{超链接}

浙江大学物理实验中心网站：\href{https://phylab.zju.edu.cn/}{https://phylab.zju.edu.cn/}

\end{document}
